\subsection{弱拓扑}

定义 2. --- 设 F 和 G 是由双线性型 B 置于对偶的两个向量空间。F 上使所有线性型 \(\mathbf{B}(., y): x \mapsto \mathbf{B}(x, y)\) （其中 y 取遍 G）都连续的最粗拓扑，称为由 F 与 G 之间的对偶定义的 F 上的弱拓扑，我们把它记作 \(\sigma(\mathbf{F}, \mathbf{G})\) 。

类似地，我们定义弱拓扑 \(\sigma(\mathbf{G},\mathbf{F})\) （\(\mathbf{G}\) 上的；在定义 1 中交换 \(\mathbf{F}\) 与 \(\mathbf{G}\) ）；这种交换 \(\mathbf{F}\) 与 \(\mathbf{G}\) 的可能性适用于本节以下的所有结果和定义。

只要不会混淆，我们用形容词「弱」和副词「弱地」表示关于弱拓扑 \(\sigma(\mathrm{F},\mathrm{G})\) 的性质。例如我们将说「弱收敛」「弱连续函数」等。

当 \(G \subset F^{*}\) 时，记号 \(\sigma(F, G)\) 总表示由对应于在 \(F \times G\) 上的限制（典范双线性型 \((x, x^{*}) \mapsto \langle x, x^{*} \rangle\) 的限制）的对偶所定义的弱拓扑。

在不对 F 和 G 附加假设的情况下，只要没有歧义，我们常用 \(\langle x, y \rangle\) 表示双线性型 B 的值 \(\mathbf{B}(x, y)\) （在点 \((x, y)\) 处的）；我们在本节余下部分采用这一约定。

带有弱拓扑 \(\sigma(F, G)\) 的向量空间 F 称为弱空间。弱拓扑 \(\sigma(F, G)\) 是局部凸的 (II, 第 26 页, prop. 4)；更精确地说，它是 \(\mathbf{R}^G\) 的乘积拓扑在线性映射 \(\phi: x \mapsto (\langle x, y \rangle)_{y \in G}\) （\(F\) 到 \(\mathbf{R}^G\) 中的）下的逆像。它由半范数 \(x \mapsto |\langle x, y \rangle|\) （其中 \(y\) 取遍 \(G\) ）所成的集定义 (II, 第 5 页)。对每个 \(\alpha > 0\) 及有限族 \((y_i)_{1 \leq i \leq n}\) （\(G\) 的点的），设 \(W(y_1, ..., y_n; \alpha)\) 是由那些 \(x \in F\) （它们满足 \(|\langle x, y_i \rangle| \leqslant \alpha\) ，\(1 \leq i \leq n\) ）组成的集；这些集（对任意的 \(\alpha, n\) 和 \(y_i\) ）构成 \(\sigma(F, G)\) 下 0 的一个基本邻域系。注意 \(W(y_1, ..., y_n; \alpha)\) 包含 \(F\) 的由方程 \(\langle x, y_i \rangle = 0\) （ \(1 \leq i \leq n\) ）定义的那个有限余维数向量子空间。

命题 2. --- 弱拓扑 \(\sigma(F, G)\) 是 Hausdorff 的，必须且只需 \(F\) 与 \(G\) 之间的对偶在 \(F\) 中分离。

这是 II, 第 3 页, prop. 2 的一个特殊情形。

命题 3. --- 设 F 和 G 是两个成对偶的实向量空间。F 上每个对 \(\sigma(F, G)\) 连续的线性型都可写成 \(x \mapsto \langle x, y \rangle\) ，其中 \(y \in G\) 。当对偶在 G 中分离时，元 \(y \in G\) 是唯一的。

因为，说线性型 \(f\) （\(\mathbf{F}\) 上的）对 \(\sigma(\mathbf{F}, \mathbf{G})\) 连续，意味着存在有限多个点 \(y_i \in \mathbf{G}\) （ \(1 \leqslant i \leqslant n\) ），使对所有 \(x\) （\(\mathbf{F}\) 中的）有 \(|f(x)| \leqslant \sup_{1 \leqslant i \leqslant n} |\langle x, y_i \rangle|\) (II, 第 6 页, prop. 5)。于是 \(n\) 个关系 \(\langle x, y_i \rangle = 0\) （ \(1 \leqslant i \leqslant n\) ）蕴含 \(f(x) = 0\) ，从而 (A, II, § 7.5, cor. 1) 存在线性组合 \(y = \sum_{i=1}^{n} \lambda_i y_i\) ，使 \(f(x) = \langle x, y \rangle\) （对所有 \(x \in \mathbf{F}\) ）。唯一性由 \((\mathrm{D}_{\mathrm{II}})\) 得出。

换言之，当对偶在 G 中分离且 F 带有拓扑 \(\sigma(\mathrm{F}, \mathrm{G})\) 时，我们可以把 G 典范地等同于对这个拓扑的 F 的对偶 (II, 第 42 页, def. 1)。

推论 1. --- F 的点组成的族 \((a_i)\) 对拓扑 \(\sigma(\mathrm{F},\mathrm{G})\) 是完全的，必须且只需对 G 中每个 \(y \neq 0\) ，存在指标 \(\iota\) 使 \(\langle a_{\iota}, y \rangle \neq 0\) 。

因为利用 prop. 3 及 I, 第 13 页, th. 1，该性质表达了这样的事实：对 \(\sigma(F, G)\) 没有闭超平面包含所有 \(a_{i}\) ；于是推论由 II, 第 38 页的 cor. 3 得出。

推论 2. --- F 的点组成的族 \((a_{\iota})\) 对拓扑 \(\sigma(\mathrm{F}, \mathrm{G})\) 是拓扑无关的，必须且只需对每个指标 \(\iota\)，存在元 \(b_{\iota} \in G\) ，使 \(\langle a_{\iota}, b_{\iota} \rangle \neq 0\) 且 \(\langle a_{\kappa}, b_{\iota} \rangle = 0\) （对所有 \(\kappa \neq \iota\) ）。

这意味着，对所有 \(\iota\) ，存在 \(\sigma(F, G)\) 中的一个闭超平面，它包含所有的 \(a_{\kappa}\) （带指标 \(\kappa \neq \iota\) 的）但不含 \(a_{\iota}\) 。

推论 3. --- 设 \(G_{1}\) 和 \(G_{2}\) 是 \(F^{*}\) 的两个向量子空间，它们（对典范双线性型的限制）与 F 成对偶。则 \(\sigma(F, G_{2})\) 比 \(\sigma(F, G_{1})\) 细，必须且只需 \(G_{1} \subset G_{2}\) 。

条件显然是充分的；反之，若 \(\sigma(F, G_2)\) 比 \(\sigma(F, G_1)\) 细，则每个对 \(\sigma(F, G_1)\) 连续的线性型也对 \(\sigma(F, G_2)\) 连续，从而由 prop. 3 得 \(G_1 \subset G_2\) 。

推论 4. --- 设 G 是向量空间 F 的对偶 F* 的一个向量子空间。则 F 与 G （对典范双线性型）成分离对偶，必须且只需 G 在 F* 中对拓扑 σ(F*, F) 稠密。

这由 cor. 1 得出。

